This section contextualizes RECTOR across five thematic axes that align with its core design decisions: multi-modal trajectory generation, constraint-enforcement strategies, rule representation and priority encoding, selection and reranking mechanisms, and system-level explainability.  Although modern predictors are increasingly accurate and diverse, the \emph{final} selected trajectory often still lacks an explicit guarantee of safety or rule compliance.  RECTOR addresses this gap by treating selection as a first-class decision stage governed by a deterministic, tiered rulebook.

%-------------------------------------------------------------------------------
\subsection{Candidate Trajectory Generation}
\label{subsec:candidate_generation}
%-------------------------------------------------------------------------------

The first axis concerns how candidate trajectories are produced, because the quality and diversity of these candidates determine the upper bound on what any downstream selector---including RECTOR---can achieve.

Most motion prediction models produce a set of candidate trajectories and then score or select them using likelihood or confidence.  A large family uses structured anchor sets, goal proposals, or lane-graph reasoning to generate multimodal futures efficiently.  Representative methods include LaneGCN~\cite{liang2020lanegcn}, TNT and DenseTNT~\cite{zhao2021tnt,gu2021densetnt}, Transformer-based set predictors such as MTR~\cite{shi2022motion} and Wayformer~\cite{nayakanti2023wayformer}, and query-centric architectures such as QCNet~\cite{zhou2023qcnet}.  Game-theoretic interaction modeling (e.g., GameFormer~\cite{huang2023gameformer}) further captures multi-agent coupling, but it still selects the final mode based on confidence.  These methods achieve strong best-of-$K$ accuracy, yet their training objectives are primarily geometric and probabilistic and therefore do not enforce traffic-rule acceptability for the executed mode.

A second family relies on latent-variable generative modeling to capture long-tail behavior and interactive uncertainty.  DESIRE~\cite{lee2017desire} and PRECOG~\cite{rhinehart2019precog} are representative early approaches, while Trajectron++~\cite{salzmann2020trajectron} models multi-agent interactions with structured latent variables.  More recent work uses large-capacity sequence models and large-scale pretraining---for example, MotionLM~\cite{seff2023motionlm}---and diffusion-based generation for flexible sampling, such as MotionDiffuser~\cite{jiang2023motiondiffuser}.  A related line of work applies constraint guidance during diffusion-based generation: CTG~\cite{zhong2023ctg} and CTG++~\cite{zhong2023ctgpp} guide trajectory sampling toward rule-satisfying regions.  \textsc{RECTOR}'s post-generation strategy is intentionally orthogonal to these constraint-guided approaches.  By decoupling selection from generation, \textsc{RECTOR} avoids both the loss of candidate diversity and the \emph{auditability gap} inherent in systems where compliance is embedded in the generation process rather than expressed as an explicit, rule-level decision trace.

Both families---anchor-based and generative---increase candidate diversity (or, in the constraint-guided case, targeted compliance), but they also amplify a practical issue: as the number of candidates grows, the decision of \emph{which} candidate to execute becomes more consequential, and confidence-only selection can be unreliable in rare scenes where the model is under-calibrated.  RECTOR is intentionally orthogonal to the generator.  It can operate on top of any multi-modal model---anchor-based, latent-variable, Transformer, or diffusion---and focuses solely on the selection step: among diverse candidates, it deterministically prioritizes those that satisfy safety and legal constraints rather than relying only on probability mass.

%-------------------------------------------------------------------------------
\subsection{When Are Rules Enforced? Constraint Entry Taxonomy}
\label{subsec:constraint_entry_taxonomy}
%-------------------------------------------------------------------------------

The previous subsection described \emph{what} candidates are generated; this subsection asks \emph{when} rules and safety constraints enter the pipeline, because the entry point determines both the guarantees a method can offer and its characteristic failure modes.  Three entry points are relevant.

\textit{Planning-time enforcement} integrates constraints directly into a control or optimization loop.  Model predictive control (MPC)~\cite{borrelli2005mpc} and control barrier functions (CBFs)~\cite{ames2017control} are common representatives, and reachability analysis provides an additional formal tool~\cite{althoff2010reachability}.  These methods can offer strong guarantees under their modeling assumptions, but their use in large-scale data-driven prediction is limited by perceptual uncertainty, interaction complexity, and the need for real-time-constrained optimization at every step.

\textit{Policy-time enforcement} introduces constraints during learning, for example via constrained reinforcement learning~\cite{altman1999constrained,achiam2017constrained} or shielding and safety filters such as SafetyNet~\cite{vitelli2022safetynet}.  These methods can empirically improve safety, but the learned constraint satisfaction may be brittle under distribution shifts, and the resulting behavior can be difficult to audit when constraints are not explicitly represented in the decision output.

\textit{Selection-time enforcement}, by contrast, evaluates candidate trajectories against rules \emph{after} generation and selects the best compliant option.  This strategy remains underexplored, yet it is natural in multi-modal settings: if the generator already provides diversity, the selector should exploit it to avoid unsafe or illegal futures.  \textsc{RECTOR} adopts this strategy, using smooth proxies and lexicographic tiering to prevent undesirable trade-offs while preserving full candidate diversity.  Critically, selection-time enforcement yields an explicit, rule-level decision trace for post hoc auditing rather than embedding compliance in opaque generation internals.

%-------------------------------------------------------------------------------
\subsection{Rule Representation and Priority Structures}
\label{subsec:rule_representation_priority}
%-------------------------------------------------------------------------------

Having established that RECTOR enforces rules at selection time, the next question is how to \emph{represent} those rules and how to encode the priority relationships among them.

In practice, industry systems often rely on structured rule checkers and safety envelopes.  RSS (Responsibility-Sensitive Safety)~\cite{shalev2017formal} is a well-known example that formalizes safe longitudinal and lateral distances together with responsibility constraints.  Its operating point is fundamentally different from RECTOR's: RSS is a planning/control-time formal safety envelope with responsibility semantics, whereas RECTOR is a selection-time ranking layer over a learned candidate set.  Censi~et~al.~\cite{censi2019liability} introduced the concept of \emph{liability-aware rulebooks} with lexicographic prioritization for autonomous driving, establishing that traffic rules should be organized into priority tiers in which higher-priority rules strictly dominate lower ones.  This rulebook concept has since influenced multiple planning frameworks that use priority-ordered constraint satisfaction.

In learning-based prediction, however, rule constraints are often reduced to scalar penalties, which introduces a familiar issue from multi-objective optimization: weight tuning implicitly defines trade-offs, and small weight changes can reorder solutions in unintended ways (see, e.g., surveys in multi-objective methods~\cite{marler2004survey}).  This is particularly problematic for traffic rules, where priorities are not subjective preferences but normative requirements---safety and legality \emph{should} dominate comfort, and that dominance should not depend on a particular weight vector.

RECTOR extends the rulebook concept of Censi~et~al.~\cite{censi2019liability} by integrating lexicographic selection into a \emph{learned} multi-modal prediction pipeline.  The specific additions beyond the original rulebook framework are threefold: (i)~a neural applicability head that predicts rule relevance from scene context, removing the need for hand-coded activation logic; (ii)~smooth differentiable proxy functions that enable gradient-based training-time shaping; and (iii)~a tolerance-based selector compatible with batched GPU inference, in which tolerances are used only to resolve near-ties so that the selector remains stable while preserving the intended hierarchy.  To the best of our knowledge, prior AV prediction and selection pipelines do not treat rule applicability itself as a learned, scene-conditioned prediction target; relevance is typically hard-coded from map logic or left implicit in downstream rule checks.

%-------------------------------------------------------------------------------
\subsection{Selection and Reranking over Candidate Sets}
\label{subsec:selection_reranking}
%-------------------------------------------------------------------------------

The preceding subsections described how candidates are generated, when constraints are applied, and how rules are prioritized.  This subsection turns to the specific mechanism of \emph{choosing} the final mode from the candidate set---the step that RECTOR targets.

Most multimodal predictors select the final mode either by taking the mode with maximum confidence (argmax over mode probabilities) or by sampling.  Although this is computationally attractive, confidence is not a reliable proxy for acceptability, especially in rare interactive scenarios where the model is under-calibrated.  An alternative is to rerank candidates using a learned scorer, a cost model, or a downstream-task objective.  In adjacent areas, differentiable optimization layers (e.g., OptNet and related differentiable solvers) have been used to impose structured constraints during learning~\cite{amos2017optnet,amos2018differentiable}.  Integrating full traffic-rule semantics into a differentiable solver remains difficult; moreover, strict constraints can collapse multimodality if they are applied too aggressively.

RECTOR addresses this gap with a non-parametric reranking mechanism that adds two components often missing from prior reranking approaches: (i)~\textit{applicability gating}, which activates only the rules relevant to the current scene and thereby avoids penalizing candidates for constraints that do not apply, and (ii)~\textit{tiered lexicographic ranking}, which prevents undesirable cross-tier trade-offs.  Together, these features keep reranking deterministic and interpretable while maintaining real-time execution and compatibility with any upstream generator.

%-------------------------------------------------------------------------------
\subsection{Explainability and Auditability}
\label{subsec:explainability_auditability}
%-------------------------------------------------------------------------------

The selection mechanism just described not only picks a trajectory but also produces a structured record of \emph{why} that trajectory was chosen---a property that connects naturally to the broader challenge of explainability in trajectory prediction.

Explainability in this domain is difficult because model decisions are distributed over time, across agents, and within latent representations.  Post hoc explainers can provide local attributions, but they may not align with the true decision boundary and are often unstable across similar inputs~\cite{ribeiro2016why,lakkaraju2016interpretable}.  In safety-critical autonomy, a more actionable form of explainability is \emph{auditability}: the system should report which constraints were considered, which were violated, and which priority level dominated the decision.

RECTOR satisfies this requirement by construction.  For each candidate, it produces tier-wise scores and rule-level proxy violations; the final decision can be summarized by the first tier that differentiates among competing candidates.  This output is closer to a safety-case artifact than to a generic feature-attribution method, and it supports debugging, dataset curation, and regression testing when rules change.

%-------------------------------------------------------------------------------
\subsection{Benchmarks and Compliance Metrics}
\label{subsec:benchmarks_metrics}
%-------------------------------------------------------------------------------

A related and complementary trend concerns the metrics and benchmarks used to evaluate trajectory predictors, because the choice of metric shapes research priorities.

Standard motion-forecasting benchmarks have historically emphasized geometric coverage of the ground truth.  Datasets such as Argoverse~\cite{chang2019argoverse} and nuScenes~\cite{caesar2020nuscenes} popularized minADE and minFDE as primary metrics, which measure whether \emph{some} predicted mode is close to the ground-truth future.  This is appropriate for measuring multi-modal coverage, but it does not assess whether the \emph{selected} mode is acceptable.  More recent planning-aware benchmarks and closed-loop evaluations---including nuPlan~\cite{caesar2021nuplan} and scenario-centric evaluations in the Waymo ecosystem~\cite{montali2024waymo}---reinforce an important point: once a model is used to select an ego trajectory, the relevant metrics must include compliance and safety, not only geometric error.

RECTOR aligns with this direction by reporting selected-trajectory metrics (selADE/selFDE) alongside explicit rule-violation rates.  Moreover, it distinguishes between proxy-based selection (for stable scoring and training compatibility) and full-catalog auditing (for semantic coverage), which matters when some rules depend on signals that are not available as smooth proxies.

%-------------------------------------------------------------------------------
\subsection{Summary: Positioning of RECTOR}
\label{subsec:related_work_summary}
%-------------------------------------------------------------------------------

Taken together, the five axes reviewed above show that prior work provides strong multimodal candidate generators, mature constraint-enforcement methods for control and planning, several rule formalization frameworks, and an evolving set of compliance-aware benchmarks.  What remains is a practical gap in \emph{selection}: given a multimodal set of candidate ego futures, how should a system choose the trajectory to execute under explicitly prioritized traffic rules while remaining real-time and auditable?

RECTOR targets this gap as a selection layer that is architecturally independent of the upstream generator.  It predicts rule applicability from scene context, evaluates tiered proxy severities, and applies tolerance-based lexicographic ordering to select a final trajectory while producing an auditable decision trace.  In this paper, we instantiate RECTOR on top of a single CVAE generator to provide an end-to-end reference stack; the selection layer is designed to be composable with other upstream models, but this composability has not been empirically validated on additional generators.  Table~\ref{tab:related_work_comparison} summarizes this positioning across representative methods from each of the five axes.

\begin{table*}[t]

\centering

\small

\caption{Positioning of RECTOR as a Generator-Agnostic Selection Layer Relative to Related Work}

\label{tab:related_work_comparison}

\begin{tabular}{l|ccccc}

\toprule

\textbf{Method} & \textbf{Multi-Modal} & \textbf{Rule-Aware (Decision)} & \textbf{Tiered Priority} & \textbf{Decision Trace} & \textbf{Real-Time} \\

\midrule

\multicolumn{6}{l}{\textit{Multi-Modal Prediction Methods}} \\

TNT/DenseTNT~\cite{zhao2021tnt,gu2021densetnt} & \checkmark & --- & --- & --- & \checkmark \\

Trajectron++~\cite{salzmann2020trajectron} & \checkmark & --- & --- & --- & \checkmark \\

MTR~\cite{shi2022motion} & \checkmark & --- & --- & --- & \checkmark \\

Wayformer~\cite{nayakanti2023wayformer} & \checkmark & --- & --- & --- & \checkmark \\

MotionDiffuser~\cite{jiang2023motiondiffuser} & \checkmark & --- & --- & --- & $\sim$ \\

QCNet~\cite{zhou2023qcnet} & \checkmark & --- & --- & --- & \checkmark \\

GameFormer~\cite{huang2023gameformer} & \checkmark & --- & --- & --- & \checkmark \\

CTG/CTG++~\cite{zhong2023ctg,zhong2023ctgpp} & \checkmark & \checkmark & --- & --- & $\sim$ \\

\midrule

\multicolumn{6}{l}{\textit{Constraint Enforcement and Safety}} \\

MPC-based~\cite{borrelli2005mpc} & --- & \checkmark & $\sim$ & \checkmark & $\sim$ \\

CBF-based~\cite{ames2017control} & --- & \checkmark & $\sim$ & \checkmark & $\sim$ \\

Constrained RL~\cite{achiam2017constrained} & --- & \checkmark & --- & --- & \checkmark \\

RSS~\cite{shalev2017formal} & --- & \checkmark & --- & \checkmark & $\sim$ \\

SafetyNet~\cite{vitelli2022safetynet} & --- & \checkmark & --- & --- & \checkmark \\

Rulebooks~\cite{censi2019liability} & --- & \checkmark & \checkmark & \checkmark & --- \\

\midrule

\multicolumn{6}{l}{\textit{RECTOR (This Work)}} \\

\textbf{RECTOR (Ours)} & \checkmark & \checkmark & \checkmark & \checkmark & \checkmark \\

\bottomrule

\end{tabular}

\vspace{2mm}



\small

\textbf{Legend:} \checkmark~= Full support. $\sim$~= Partial/conditional support. ---~= Not supported.

\vspace{1mm}



\small

\textbf{Note:} Rows under \textit{Multi-Modal Prediction Methods} are \emph{candidate generators}.  They do not natively include rule-aware, tiered decision logic; attaching RECTOR to any such generator yields a composite stack that is rule-aware, hierarchical, and interpretable \emph{at selection time}.  Characterisations in this table are based on published descriptions of each method as research baselines and may not reflect the current state of any commercial implementation.

\end{table*}